% Corde a t = T/8 : translation de lambda/8 vers les x croissants
% \begin{tikzpicture}[scale=1]
%     \def\lam{4}
%     \def\amp{1}
%     \def\shift{0.5} % lambda/8
%     \draw[->] (-0.3,0) -- (8.5,0) node[right] {$x$};
%     \draw[->] (0,-1.5) -- (0,1.5) node[above] {$y$};
%     \draw[thick, gray!50, domain=0:8, samples=200, smooth]
%         plot (\x, {\amp*sin(2*180/\lam*\x)});
%     \draw[thick, red, domain=0:8, samples=200, smooth]
%         plot (\x, {\amp*sin(2*180/\lam*(\x-\shift))});
% \end{tikzpicture}

% \begin{tikzpicture}[scale=1]
%     \def\lam{4}
%     \def\amp{1.2}
%     \def\shift{0.5} % lambda/8

%     % 1. Grille stylisée (arrière-plan)
%     \drawgrid{-1}{9}{-2}{2}

%     % 2. Axes au style "Mer de Fermi"
%     \draw[fermi axis] (-0.5, 0) -- (8.8, 0) node[pos=1.02, below] {\huge $x$};
%     \draw[fermi axis] (0, -1.8) -- (0, 1.8) node[pos=1.02, left] {\huge $y$};

%     % 3. Courbe 1 (Signal original) -> Bleu Indigo (colorFour)
%     \draw[wave primary, domain=0:8, samples=200, smooth]
%         plot (\x, {\amp*sin(2*180/\lam*\x)});

%     % 4. Courbe 2 (Signal déphasé) -> Rouge Corail (colorSix)
%     \draw[wave secondary, domain=0:8, samples=200, smooth]
%         plot (\x, {\amp*sin(2*180/\lam*(\x-\shift))});

%     % 5. Légende / Annotations harmonisées
%     \node[color=colorFour, right] at (4.5, 1.4) {$\psi_1(x)$};
%     \node[color=colorSix, right] at (6.0, -1.4) {$\psi_2(x - \delta)$};

% \end{tikzpicture}

% Figure 1 -- Allure de la corde a t=0 (onde sinusoidale, lambda = 1 unite)
% \begin{tikzpicture}[scale=1]
%     \def\lam{4}   % longueur d'onde en cm sur la figure
%     \def\amp{1}   % amplitude en cm sur la figure
%     \draw[->] (-0.3,0) -- (8.5,0) node[right] {$x$};
%     \draw[->] (0,-1.5) -- (0,1.5) node[above] {$y$};
%     \draw[thick, blue, domain=0:8, samples=200, smooth]
%         plot (\x, {\amp*sin(2*180/\lam*\x)});
%     \draw[dashed] (\lam,-1.3) -- (\lam,1.3);
%     \draw[<->] (0,-1.3) -- (\lam,-1.3) node[midway, below] {$\lambda$};
%     \node[blue] at (1,1.3) {$P$};
%     \fill[blue] (1,{\amp*sin(2*180/\lam)}) circle (1.5pt);
% \end{tikzpicture}


\begin{tikzpicture}
	

	% Définition de l’échelle
	%\def\echelleY{10}
	\pgfmathsetmacro{\echelleY}{2}
	%\def\echelleX{50}
	\pgfmathsetmacro{\echelleX}{1.5}

    %\clip[scale = 0.95 , decorate, decoration={random steps, segment length=0.1 cm, amplitude=0.2cm}](#1 , #2 ) rectangle ( #3 , #4 ) ;
	\clip[](-1.4*\echelleX , -2*\echelleY) rectangle (10.5*\echelleX , 2.0*\echelleY) ;

	% Définition de la liste
	\def\listeLabelsY{
        -1/-1/,
        0//,
        1/1/
    }

    \def\listeLabelsX{
        -1/-1/,
        0//,
        1/1/,
        2/2/,
        3/3/,
        4/4/,
        5/5/,
        6/6/,
        7/7/,
        8/8/,
        9/9/,
        10/10/
    }

    \def\lam{4}   % longueur d'onde en cm sur la figure
    \def\amp{1}   % amplitude en cm sur la figure
    \def\phi{-2*180/\lam*0.5}  % phase en cm sur la figure 

	\begin{scope}

		%\clip[](\valXmin , -1) rectangle (9, 11.0) ;
        
		\backgrowngrid{-1}{-3}{10}{3}{\echelleX}{\echelleY/2}{\colorThree!50!\colorTwo}{\colorTwo}{0.5ex}
        

        \begin{scope}[shift={(1*\echelleX,1.25*\echelleY)},]
            \draw[
                line width=0.5ex,
                draw=\colorThree,
                fill=\colorThree!50!\colorTwo,
                rounded corners=6pt,
                opacity = 0.5
            ]
            (0,0) rectangle (9*\echelleX,0.5*\echelleY);


            \draw[shift={(0.5*\echelleX,0.25*\echelleY)}]
                (0.0*\echelleX,0*\echelleY) edge[shift = {(0,0)},line width=0.8ex  ,color=\colorFour] node[shift = {(0,0)} , pos=1, right , font=\fontsize{12pt}{12pt}\selectfont \bf, text=\colorFour 
                    ]{${\bm f(x) = y(x , t = 0 )}$} (1.0*\echelleX,0.0*\echelleY)

                (0.0*\echelleX,0*\echelleY) edge[shift = {(6,0)} , line width=0.8ex  ,color=\colorSix] node[shift = {(0,0)} , pos=1, right , font=\fontsize{12pt}{12pt}\selectfont \bf, text=\colorSix 
                    ]{${\bm f(x) = y(x , t = T/8 )}$} (1.0*\echelleX,0.0*\echelleY)
                ;
        \end{scope}
        
        \draw[opacity = 1, line width=0.8ex  ,color=colorFour ,  fill = none, domain=-2:11, samples=200, smooth]
            plot (\x*\echelleX, {\echelleY*\amp*cos(2*180/\lam*\x + \phi)}) 
            ;

        \draw[line width=0.8ex  ,color=\colorSix ,  fill = none, domain=-2:11, samples=200, smooth]
            plot (\x*\echelleX, {\echelleY*\amp*cos(2*180/\lam*(\x-\lam/8) + \phi)}) 
            ;

        
		% \draw[line width=0.8ex  ,color=\colorFour ,  fill = none]
		% plot[smooth] file {data/Donnee_article_simul_expansion_1_24-04-2024-T-x0.table}
		% (-10.2,10) edge[] node[shift = {(0,0)} , pos=1, right ]{\fontsize{16pt}{19pt}\selectfont \bf Données } (-8.7,10)
		% ;
		% \draw[line width=0.8ex  ,color=\colorSix,  fill = none]
		% plot[smooth] file {data/GHD_1_article_simul_expansion_1_24-04-2024-T-x0.table}
		% (-10.2,9) edge[] node[shift = {(0,0)} , pos=1, right ]{\fontsize{16pt}{19pt}\selectfont \bf T = 560 nK  } (-8.7,9)
		% ;
		% \draw[line width=0.8ex  ,color=\colorThree,  fill = none]
		% plot[] file {data/GHD_2_article_simul_expansion_1_24-04-2024-T-x0.table}
		% (-10.2,8) edge[] node[shift = {(0,0)} , pos=1, right ]{\fontsize{16pt}{19pt}\selectfont \bf Ajust : T = 1550 nK  } (-8.7,8)
		% ;

	\end{scope}

	% Graduations     
	\begin{scope}[shift = {(0,0)}]
		\GraduationsY{\listeLabelsY}{\echelleY}{0.3}{0.5ex}{\colorOne}{\Large}
        \GraduationsY{-1.5/-1.5/, -0.5/-0.5/ , 0.5/0.5/ , 1.5/1.5/}{\echelleY}{0.2}{0.4ex}{\colorOne}{\large}
	\end{scope}
	\begin{scope}[shift = {(0,0)}]
		\GraduationsX{\listeLabelsX}{\echelleX}{0.3}{0.5ex}{\colorOne}{\Large}
    \end{scope}

	% Axxes 
	\draw[shift = {(0,0)}]
        (-2*\echelleX, 0) 
            edge[
                thick, 
                line width=0.8ex, 
                ->, 
                >=Triangle, 
                color=\colorOne
                ] 
            node[
                shift = {(0,-1)} ,
                pos=0.85,
                below,
                font=\sffamily\bfseries\fontsize{15pt}{1pt}\selectfont\bfseries,
                text=\colorOne
                ]
            { 
                ${\bm x}~(\text{cm}) $ 
            } 
        (10.5*\echelleX, 0)
        ;
	\draw [shift = {(0,0)}]  
        (0, -1.75*\echelleY) 
            edge[
                thick, 
                line width=0.8ex, 
                ->, 
                >=Triangle, 
                color=\colorOne
                ] 
            node[
                shift = {(-1,0)} , 
                pos=0.85, 
                above , 
                rotate = 90,
                font=\sffamily\bfseries\fontsize{15pt}{1pt}\selectfont\bfseries,
                text=\colorOne ,
                ]
            {
                ${\bm y}~(\text{mm}) $ 
            } 
        (0, 1.85*\echelleY)

    ;

    \draw[
        dashed,
        line width=0.6ex,
        shift = {(1.5*\echelleX,0)} 
        ] 
        (0,-1.55*\echelleY) -- (0,0.55*\echelleY)
        (\lam/8*\echelleX,-1.55*\echelleY) -- (\lam/8*\echelleX,0.55*\echelleY)
        ;
    \draw[
        thick, 
        line width=0.8ex, 
        >-< , 
        shift = {(1.5*\echelleX,-1.55*\echelleY)} 
        ] 
        (0,0) --++ (\lam/8*\echelleX,0*\echelleY) 
        node[
            shift ={(0,-0.2)},
            midway, 
            below,
            font=\sffamily\bfseries\fontsize{15pt}{1pt}\selectfont\bfseries,
            text=\colorOne ,
            ] {${\bm\lambda/8}$};

    \pgfmathsetmacro{\xP}{1.5*\echelleX}
    \pgfmathsetmacro{\yP}{
        \echelleY*\amp*cos(360*1.5/\lam + \phi)
    }

    \node[
        shift = {(0.4,0.4)},
        font=\sffamily\bfseries\fontsize{20pt}{1pt}\selectfont\bfseries, 
        text=\colorFour
    ]
    at (\xP,\yP)
    {$P$};

    \fill[
        color=\colorFour
    ]
    (\xP,\yP)
    circle (0.15cm);
	%\draw[line width=2.5ex , color = red ] (-14*\echelleX , -1*\echelleY) rectangle (9.5*\echelleX , 12*\echelleY) ;
	%\drawgrid{-2}{11}{-2}{2}{\echelleX}{\echelleY}

\end{tikzpicture}

